\documentclass{article}
\usepackage{graphicx} % Required for inserting images
\usepackage{amsmath}
\usepackage{amssymb}
\usepackage{amsthm}
\usepackage{float}

\title{MATH 402 Homework 6}
\author{Aaron Honculada}
\date{February 2025}

\begin{document}

\maketitle

\section{Problem 1}
\subsection{Show that $\sqrt{p}$ is irrational for every positive prime integer $p$.}

\begin{proof}
    We use proof by contradiction. Suppose that $\sqrt{p}$ is rational. Then we can write $\sqrt{p} = \frac{a}{b}$ where $a$ and $b$ are integers.
    Without loss of generality, we can assume that $a$ and $b$ are coprime, i.e. $\gcd(a,b) = 1$. Then we have $p = \frac{a^2}{b^2}$ which implies that $p \mid a^2$. Since $p$ is prime,
    we must have $p \mid a$. Therefore, we can write $a = kp$ for some integer $k$. Substituting, we get:
    \[
        p = \frac{k^2p^2}{b^2} \implies p(b^2) = k^2p^2 \implies b^2 = k^2p \implies p \mid b^2.
    \]
    Since $p$ is prime, we must have $p \mid b$. But this contradicts the fact that $\gcd(a,b) = 1$ since it was shown that $p \mid a$ and $p \mid b$. Therefore, $\sqrt{p}$ must be irrational
    for every positive prime integer $p$.
\end{proof}

\section{Problem 2}
\subsection{Show there are infinitely many integers $k$ such that $x^9 + 12x^5 + 21x + k$ is irreducible over $\mathbb{Q}[x]$.}

\begin{proof}
    We use Eisenstein's criterion to show that $x^9 + 12x^5 + 21x + k$ is irreducible over $\mathbb{Q}[x]$ for some prime $p$ if:
    \begin{itemize}
        \item $p \mid 12$
        \item $p \mid 21$
        \item $p \nmid 1$
        \item $p^2 \nmid k$
    \end{itemize}
    We choose $p=3$ and the first three conditions are satisfied. Recall that the set of integers satisfying $k \equiv 0 \pmod{3}$ but 
    $k \not\equiv 0 \pmod{9}$ is infinite. This is the set $3\mathbb{Z} \setminus 9\mathbb{Z}$ whose cardinality is countably infinite. Therefore
    there are infinitely many integers $k$ such that $x^9 + 12x^5 + 21x + k$ is irreducible over $\mathbb{Q}[x]$.
\end{proof}

\section{Problem 3}
\subsection{Prove that for $p$ prime, $f(x) = x^{p-1} + x^{p-2} + \cdots + x + 1$ is irreducible in $\mathbb{Q}[x]$.}

\begin{proof}
    We consider $f(x)$ as the sum of the first $p$ terms of a geometric series:
    \[
        f(x) = \frac{x^p - 1}{x-1} \text{ for } x \neq 1.
    \]
    Now we make a change of variable $x = y + 1$ to get:
    \[
        g(y) = \frac{(y+1)^p - 1}{y}.
    \]
    If we show that $g(y)$ is irreducible over $\mathbb{Q}[y]$, then $f(x)$ is irreducible over $\mathbb{Q}[x]$ since $x = y+1$.
    Using the binomial theorem:
    \[
        (y+1)^p - 1 = y^p + \binom{p}{1}y^{p-1} + \cdots + \binom{p}{p-1}y + 1.
    \]
    Therefore, we can write:
    \[
        g(y) = \frac{\sum_{k=0}^{p} \binom{p}{k}y^k-1}{y} = \frac{1 + py + \binom{p}{2}y^2 + \cdots + y^p - 1}{y}.
    \]
    Simplifying the numerator, we get:
    \[
        g(y) = p + \binom{p}{2}y + \binom{p}{3}y^2 + \cdots + y^{p-1}.
    \]
    We use Eisenstein's criterion for any prime $p$ to show that $g(y)$ is irreducible over $\mathbb{Q}[y]$.
    \begin{itemize}
        \item $p$ does not divide the leading coefficient, which is $1$.
        \item $p$ divides each coefficient of $g(y)$ except for the leading coefficient.
        \item $p^2$ does not divide the constant term, which is $p$.
    \end{itemize}
    All of the conditions of Eisenstein's criterion are satisfied, therefore $g(y)$ is irreducible over $\mathbb{Q}[y]$.
    Therefore, $f(x)$ is irreducible over $\mathbb{Q}[x]$.
\end{proof}

\newpage
\section{Problem 4}
Factor each polynomial as a product of irreducible polynomials in $\mathbb{Q}[x]$, $\mathbb{R}[x]$, and $\mathbb{C}[x]$.
\subsection{Solution}
\begin{table}[h]
    \centering
    \begin{tabular}{c|p{6cm}}
        $p(x)$ & $\mathbb{Q}[x]$ \\
        \hline
        $x^4-2$ & $(x^2-\sqrt{2})(x^2+\sqrt{2})$ \\
        $x^3+1$ & $(x+1)(x^2-x+1)$ \\
        $x^3-x^2-5x+5$ & $(x-1)(x^2-5)$ \\
    \end{tabular}
    \caption{Factorization in $\mathbb{Q}[x]$}
\end{table}

\begin{table}[h]
    \centering
    \begin{tabular}{c|p{6cm}}
        $p(x)$ & $\mathbb{R}[x]$ \\
        \hline
        $x^4-2$ & $(x-\sqrt[4]{2})(x+\sqrt[4]{2})(x^2+\sqrt{2})$ \\
        $x^3+1$ & $(x+1)(x^2-x+1)$ \\
        $x^3-x^2-5x+5$ & $(x-\sqrt{5})(x+\sqrt{5})(x-1)$ \\
    \end{tabular}
    \caption{Factorization in $\mathbb{R}[x]$}
\end{table}

\begin{table}[h]
    \centering
    \begin{tabular}{c|p{6cm}}
        $p(x)$ & $\mathbb{C}[x]$ \\
        \hline
        $x^4-2$ & $(x-\sqrt[4]{2})(x+\sqrt[4]{2})(x-i\sqrt[4]{2})(x+i\sqrt[4]{2})$ \\
        $x^3+1$ & $(x+1)(x-\frac{1+i\sqrt{3}}{2})(x-\frac{1-i\sqrt{3}}{2})$ \\
        $x^3-x^2-5x+5$ & $(x-\sqrt{5})(x+\sqrt{5})(x-1)$ \\
    \end{tabular}
    \caption{Factorization in $\mathbb{C}[x]$}
\end{table}

\section{Problem 5}
\subsection{Factor $x^2+x+1+i$ in $\mathbb{C}[x]$.}

\begin{proof}
    Using the quadratic formula, the roots are given by:
    \[
        x = \frac{-b \pm \sqrt{b^2-4ac}}{2a}
    \]
    Therefore, we can factor $x^2+x+1+i$ as:
    \[
        \left(x-\frac{-1+(1-2i)}{2}\right)\left(x-\frac{-1-(1-2i)}{2}\right)
    \]
\end{proof}

\section{Problem 6}
\subsection{Show that a polynomial of odd degree in $\mathbb{R}[x]$ with no multiple roots must have an odd number of real roots.}

\begin{proof}
    Let $f(x) = a_n x^n + a_{n-1} x^{n-1} + \cdots + a_1 x + a_0$ be a polynomial of odd degree $n$ with no multiple roots.
    We can write $f(x) = (x-r_1)(x-r_2)\cdots(x-r_n)$ where $r_1, r_2, \ldots, r_n$ are the roots of $f(x)$.
    Since $f(x)$ has no multiple roots, all the $r_i$ are distinct. Since $f(x)$ has odd degree, at least one of the $r_i$ must be real.
    This follows from the fact that $f(x)$ has real coefficients, the complex roots must come in conjugate pairs. Therefore, the number of real roots must be odd.
    That is to say that if $a+bi$ is a root, then $a-bi$ is also a root. Thus, the number of complex roots is always even. If $k$ complex roots exist, they must come
    in $\frac{k}{2}$ conjugate pairs.

    Let $r$ be the number of real roots. Then the number of complex roots is $n-r$ which must be even since they come in conjugate pairs. This implies that:
    \[
        n - r \text{ is even.}
    \]
    We rewrite this as:
    \[
        r = n - (2k \text{ for some integer } k) = odd - even = odd.
    \]
    Therefore, the number of real roots must be odd.
\end{proof}

\section{Problem 7}
Let $f(x), g(x), p(x) \in F[x]$, with $p(x)$ non-zero. Determine whether $f(x) \equiv g(x) \pmod{p(x)}$.
\subsection{$f(x) = x^5 -2x^4 + 4x^3 + x + 1, g(x) = 3x^4 + 2x^3 - 5x^2 - 9, p(x) = x^2 + 1, F = \mathbb{Q}$}
We reduce $f(x)$ and $g(x)$ modulo $p(x)$, which is essentially replacing $x^2$ with $-1$:
\begin{itemize}
    \item $x^5 = x(x^4) = x(x^2)^2 = x((-1)^2) = x$
    \item $x^4 = x^2 \cdot x^2 = (-1)^2 = 1$
    \item $x^3 = x \cdot x^2 = x(-1) = -x$
    \item $x^2 = -1$
\end{itemize}

\[
    f(x) \equiv x^5 -2x^4 + 4x^3 + x + 1 \equiv -2x - 1 \pmod{x^2 + 1}
\]
\[
    g(x) \equiv 3x^4 + 2x^3 - 5x^2 - 9 \equiv -2x - 1 \pmod{x^2 + 1}
\]

Therefore we have $f(x) \equiv g(x) \pmod{x^2 + 1}$.

\subsection{$f(x) = x^4 + x^2 + x + 1, g(x) = x^4 + x^3 + x^2 + 1, p(x) = x^2 + x, F = \mathbb{Z}_2$}
We find the difference between $f(x)$ and $g(x)$:
\[
    f(x) - g(x) = (x^4 + x^2 + x + 1) - (x^4 + x^3 + x^2 + 1) = x^3 + x
\]
We reduce $x^3 + x$ modulo $x^2 + x$ to get:
\[
    x^3 + x \equiv x \pmod{x^2 + x}
\]
Subtracting $x(x^2 + x)$ from both sides, we get:
\[
    x^3 + x - x(x^2 + x) = x-x^2 = x^2 + x \text{ in } \mathbb{Z}_2[x]
\]
We see that $f(x) - g(x) \equiv 0 \pmod{x^2 + x}$, therefore $f(x) \equiv g(x) \pmod{x^2 + x}$.


\subsection{$f(x) = x^5 -2x^4 + 4x^3 -4x^2 + 5x -1, g(x) = x^4 -3x^3 + 5x^2 -7x + 9, p(x) = x^2 + 1, F = \mathbb{Q}$}
By performing polynomial long division, we get:
\[
    f(x) \equiv 2x + 1 \pmod{x^2 + 1}
\]
\[
    g(x) \equiv -4x + 5 \pmod{x^2 + 1}
\]
By taking the difference, we get:
\[
    f(x) - g(x) \equiv 2x + 1 - (-4x + 5) = 6x - 4 \not\equiv 0 \pmod{x^2 + 1}
\]
We see that $f(x) - g(x) \not\equiv 0 \pmod{x^2 + 1}$, and therefore $f(x) \not\equiv g(x) \pmod{x^2 + 1}$.


\section{Problem 8}
List all the possible congruence classes in \(\mathbb{Z}_2 / (x^3 + x + 1)\).

The possible congruence classes are represented by the polynomials of degree less than 3 with coefficients in \(\mathbb{Z}_2\). These are:

\[
\begin{aligned}
    &0, \\
    &1, \\
    &x, \\
    &x+1, \\
    &x^2, \\
    &x^2+1, \\
    &x^2+x, \\
    &x^2+x+1.
\end{aligned}
\]

\section{Problem 9}
\subsection{Prove or disprove: If $p(x)$ is irreducible in $F[x]$ and $f(x)g(x) \equiv 0_F \pmod{p(x)}$, then $f(x) \equiv 0_F \pmod{p(x)}$ or $g(x) \equiv 0_F \pmod{p(x)}$.}

\begin{proof}
The congruence relation means that $p(x) \mid f(x)g(x)$. Since $p(x)$ is irreducible, it must divide either $f(x)$ or $g(x)$. This follows from
    the fact that $F[x]/p(x)$ is an integral domain which does not have zero divisors.
In other words, $f(x) \equiv 0_F \pmod{p(x)}$ or $g(x) \equiv 0_F \pmod{p(x)}$.
    In other words, $f(x) \equiv 0_F \pmod{p(x)}$ or $g(x) \equiv 0_F \pmod{p(x)}$.
\end{proof}



\section{Problem 10}
\subsection{If $f(x)$ is relatively prime to $p(x)$, prove that there is a polynomial $g(x) \in F[x]$ such that $f(x)g(x) \equiv 1_F \pmod{p(x)}$.}

\begin{proof}
Since $f(x)$ and $p(x)$ are relatively prime, we know from Bezout's identity that there exist polynomials $a(x)$ and $b(x)$ such that $a(x)f(x) + b(x)p(x) = 1$.
Rearranging, we get $a(x)f(x) = 1 - b(x)p(x)$. And taking the congruence of both sides modulo $p(x)$, we get $a(x)f(x) \equiv 1 \pmod{p(x)}$.
This shows that $a(x)$ is the inverse of $f(x)$ in the quotient ring $F[x] / p(x)$.

Thus, for any polynomial $f(x) \in F[x]$ that is relatively prime to $p(x)$, there exists a polynomial $g(x) \in F[x]$ such that $f(x)g(x) \equiv 1_F \pmod{p(x)}$.
\end{proof}


\section{Problem 11}
Write out the addition and multiplication tables for the congruence class rings $F[x]/p(x)$. In each case, is $F[x]/p(x)$ a field?

\subsection{Let $F = \mathbb{Z}_3$ and $p(x) = x^2 + 1$}

\begin{table}[H]
    \centering
    \begin{tabular}{c|ccccccccc}
        + & 0 & 1 & 2 & x & x+1 & x+2 & 2x & 2x+1 & 2x+2 \\
        \hline
        0 & 0 & 1 & 2 & x & x+1 & x+2 & 2x & 2x+1 & 2x+2 \\
        1 & 1 & 2 & 0 & x+1 & x+2 & x & 2x+1 & 2x+2 & 0 \\
        2 & 2 & 0 & 1 & x+2 & x & 2x+2 & 1 & 2x & 0 \\
        x & x & x+1 & x+2 & 2x & 2x+1 & 2x+2 & 1 & 2 & 0 \\
        x+1 & x+1 & x+2 & 0 & 2x+1 & 2x+2 & 0 & 1 & 2 & x \\
        x+2 & x+2 & 0 & 1 & 2x+2 & 1 & 2x & 0 & x & 2x+1 \\
        2x & 2x & 2x+1 & 2x+2 & 1 & 2 & 0 & x & x+1 & x+2 \\
        2x+1 & 2x+1 & 2x+2 & 1 & 2x+1 & 0 & 1 & x+1 & x+2 & 0 \\
        2x+2 & 2x+2 & 1 & 2x & 2x+2 & x & 1 & x+2 & 0 & x+1 \\
    \end{tabular}
    \caption{Addition table for \(\mathbb{Z}_3[x] / (x^2+1)\)}
\end{table}

\begin{table}[H]
    \centering
    \begin{tabular}{c|ccccccccc}
        $\cdot$ & 0 & 1 & 2 & x & x+1 & x+2 & 2x & 2x+1 & 2x+2 \\
        \hline
        0 & 0 & 0 & 0 & 0 & 0 & 0 & 0 & 0 & 0 \\
        1 & 0 & 1 & 2 & x & x+1 & x+2 & 2x & 2x+1 & 2x+2 \\
        2 & 0 & 2 & 1 & x+1 & 2x+1 & 2x+2 & 0 & 2x & 2x+1 \\
        x & 0 & x & x+1 & 2x & 1 & 2x+1 & 2x+2 & 2 & 0 \\
        x+1 & 0 & x+1 & x+2 & 2x+1 & 2x+2 & 0 & 1 & 2 & x \\
        x+2 & 0 & x+2 & 0 & 2x+2 & 1 & 2x & 2 & x & 2x+1 \\
        2x & 0 & 2x & 2x+1 & 1 & 2 & 0 & x & x+1 & x+2 \\
        2x+1 & 0 & 2x+1 & 2x+2 & 2x+1 & 0 & 1 & x+1 & x+2 & 0 \\
        2x+2 & 0 & 2x+2 & 1 & 2x+2 & x & 2 & x+2 & 0 & x+1 \\
    \end{tabular}
    \caption{Multiplication table for \(\mathbb{Z}_3[x] / (x^2+1)\)}
\end{table}

We check if $x^2 + 1$ is irreducible over $\mathbb{Z}_3$.

\begin{itemize}
    \item $x=0$: $x^2 + 1 = 1 \neq 0$
    \item $x=1$: $x^2 + 1 = 2 \neq 0$
    \item $x=2$: $x^2 + 1 = 5 \equiv 2 \pmod{3} \neq 0$
\end{itemize}
Since $x^2 + 1$ has no roots in $\mathbb{Z}_3$ it is irreducible over $\mathbb{Z}_3$, and therfore
we conclude that $\mathbb{Z}_3[x] / (x^2+1)$ is a field.


\subsection{Let $F = \mathbb{Z}_2$ and $p(x) = x^2 + 1$}

\begin{table}[H]
    \centering
    \begin{tabular}{c|ccc}
        + & 0 & 1 & x \\
        \hline
        0 & 0 & 1 & x \\
        1 & 1 & 0 & x+1 \\
        x & x & x+1 & 0 \\
    \end{tabular}
    \caption{Addition table for \(\mathbb{Z}_2[x] / (x^2+1)\)}
\end{table}
We check if $x^2 + 1$ is irreducible over $\mathbb{Z}_2$.

\begin{table}[H]
    \centering
    \begin{tabular}{c|ccc}
        $\cdot$ & 0 & 1 & x \\
        \hline
        0 & 0 & 0 & 0 \\
        1 & 0 & 1 & x \\
        x & 0 & x & 1 \\
    \end{tabular}
    \caption{Multiplication table for \(\mathbb{Z}_2[x] / (x^2+1)\)}
\end{table}

\begin{itemize}
    \item $x=0$: $x^2 + 1 = 1 \neq 0$
    \item $x=1$: $x^2 + 1 = 2 \equiv 0 \pmod{2}$
\end{itemize}
Since $x^2 + 1$ has a root in $\mathbb{Z}_2$ it is reducible over $\mathbb{Z}_2$, and therefore
we conclude that $\mathbb{Z}_2[x] / (x^2+1)$ is not a field.

\end{document}